\documentclass[11pt]{article}

\usepackage[margin=1in]{geometry}
\usepackage{booktabs}
\usepackage{hyperref}
\usepackage{url}

\title{Agent Anvil: Trace-Centric CI Evaluation for Tool-Using AI Agents}
\author{Agent Axiom}
\date{}

\begin{document}
\maketitle

\begin{abstract}
Tool-using AI agents fail in ways that are often invisible to final-answer
evaluation: an agent may produce a plausible response while calling the wrong
tool, passing hallucinated arguments, retrying unsafe operations, or invoking a
destructive tool before verification. We present Agent Anvil, a local,
CI-first evaluation harness for trace-centric agent regression testing. Agent
Anvil runs YAML scenario suites, captures model and tool traces, applies
deterministic assertions and policy checks, optionally uses OpenAI structured
semantic grading, clusters failures, and emits Markdown and JSON artifacts for
continuous integration. In a compact offline benchmark of 100 trials across
refund, tool-safety, external-agent protocol, account-admin, and data-pipeline
scenarios, a final-answer
baseline marks 100.0\% of trials as passing (95\% Wilson CI:
96.3--100.0), while Agent Anvil's trace-aware checks mark 30.0\% as
passing (95\% Wilson CI: 21.9--39.6) and expose 70 failures that answer-only
evaluation misses. These results are not a universal leaderboard; they are a
reproducible systems case study showing why agent evaluation should inspect the
behavioral path, not only the final response.
\end{abstract}

\section{Introduction}

Large language model agents increasingly call tools that read data, mutate
state, or trigger real-world side effects. Prior work has shown the importance
of interleaving reasoning and action in language-model agents
\cite{yao2023react} and of teaching models when and how to call tools
\cite{schick2023toolformer}. In production settings, however, the evaluation
problem shifts from ``did the answer sound correct?'' to ``did the agent behave
safely while getting there?''

Answer-only evaluation can miss workflow failures. A refund agent can issue a
refund with a synthetic order identifier and then respond fluently. An
operations agent can retry a locked job or scale a service before verifying its
identity. These are trace-level bugs: they live in the sequence of model calls,
tool calls, arguments, and tool results.

Agent Anvil targets this gap. It is not a hosted observability platform and not
a replacement for human domain review. It is a small, local harness intended to
turn tool-use behavior into regression tests that run in developer workflows and
CI.

\section{Design Goals}

Agent Anvil is built around five design goals.

\textbf{Trace-first evaluation.} Scenarios are evaluated against persisted
traces containing model calls, tool calls, tool arguments, tool results,
protocol errors, and final output.

\textbf{Deterministic invariants before model judgment.} Exact safety and
business invariants are encoded as deterministic checks: expected tools,
forbidden tools, required arguments, ordered calls, maximum tool-call counts,
forbidden argument values, tool-result checks, clarification behavior, and
policy preconditions.

\textbf{Optional semantic grading.} When OpenAI credentials are available,
Agent Anvil can apply structured semantic grading to criteria that are less
suited to exact checks, such as whether a clarification is meaningful or
whether a repair suggestion is appropriate.

\textbf{CI-compatible artifacts.} Runs write JSON traces, JSON results,
Markdown reports, failure clusters, and repair plans. The CLI returns a
non-zero exit code when an eval fails.

\textbf{Bring-your-own-agent protocol.} Agents can be imported as Python
modules or run as external commands that exchange a JSON scenario payload and
JSONL trace events. This keeps the harness independent of a single agent
framework.

\section{System Overview}

Agent Anvil takes a YAML scenario suite as input. A runner executes the target
agent for one or more trials per scenario. The trace recorder persists a
versioned trace artifact with schema version \texttt{anvil.trace.v1}. The
deterministic grader evaluates trace completion, tool choice, tool arguments,
assertions, and policy checks. The optional OpenAI grader evaluates semantic
criteria and emits a structured grade with failure type, severity, reason, and
repair suggestions. The report layer clusters failures and writes Markdown and
JSON artifacts.

\section{Benchmark}

We include a compact offline benchmark under \texttt{experiments/paper.yaml}.
The benchmark compares a weak final-answer baseline against Agent Anvil's
trace-aware checks over the same traces. The baseline passes a trial when the
final response exists and does not contain obvious runtime failure terms. This
baseline is intentionally simple: it represents a common insufficient
evaluation pattern, not the strongest possible answer-only judge.

\begin{table}[h]
\centering
\begin{tabular}{lrrr}
\toprule
Suite & Trials & Final-answer pass (95\% CI) & Trace-aware pass (95\% CI) \\
\midrule
Refund trace suite & 20 & 100.0\% [83.9, 100.0] & 50.0\% [29.9, 70.1] \\
Tool-safety trace suite & 30 & 100.0\% [88.6, 100.0] & 0.0\% [0.0, 11.4] \\
External protocol trace suite & 10 & 100.0\% [72.2, 100.0] & 100.0\% [72.2, 100.0] \\
Account-admin trace suite & 20 & 100.0\% [83.9, 100.0] & 0.0\% [0.0, 16.1] \\
Data-pipeline trace suite & 20 & 100.0\% [83.9, 100.0] & 50.0\% [29.9, 70.1] \\
\midrule
Total & 100 & 100.0\% [96.3, 100.0] & 30.0\% [21.9, 39.6] \\
\bottomrule
\end{tabular}
\caption{Offline benchmark results generated by \texttt{uv run anvil paper reproduce}. Intervals are 95\% Wilson confidence intervals.}
\label{tab:benchmark}
\end{table}

Across 100 trials, the final-answer baseline marks every trial as passing.
Agent Anvil's trace-aware checks expose 70 answer-only missed failures, all
classified as policy violations in the current benchmark artifact. The result
demonstrates the intended gap: final text can look acceptable while the agent's
intermediate behavior violates tool-use safety contracts.

\begin{table}[h]
\centering
\begin{tabular}{lrrr}
\toprule
Evaluator & Trials & Pass rate & Answer-only missed failures \\
\midrule
Final answer baseline & 100 & 100.0\% [96.3, 100.0] & 0 \\
Trace completion only & 100 & 100.0\% [96.3, 100.0] & 0 \\
Deterministic assertions & 100 & 30.0\% [21.9, 39.6] & 70 \\
Policy checks & 100 & 30.0\% [21.9, 39.6] & 70 \\
Full trace-aware & 100 & 30.0\% [21.9, 39.6] & 70 \\
\bottomrule
\end{tabular}
\caption{Evaluator-layer ablation over the same 100 benchmark traces. The current benchmark focuses on policy-backed unsafe tool behavior, so deterministic assertions, policy checks, and the full trace-aware evaluator expose the same 70 answer-only missed failures.}
\label{tab:ablation}
\end{table}

The ablation in Table~\ref{tab:ablation} clarifies which evaluation layers
matter for this artifact. Merely requiring a completed trace does not improve
over final-answer checking, because all benchmark agents produce completed
traces and plausible final text. The missed failures appear only when the
evaluator inspects tool arguments, ordering, forbidden calls, and risky-tool
preconditions.

\section{Artifact}

The artifact is public and reproducible. From a clean checkout:

\begin{verbatim}
uv sync --group dev
uv run anvil paper reproduce
\end{verbatim}

The command writes \texttt{docs/paper/results.json},
\texttt{docs/paper/results.md}, CSV/LaTeX table artifacts under
\texttt{docs/paper/tables/}, and local trace artifacts under
\texttt{runs/paper-benchmark/}. The generated benchmark summary reports 100
total trials, 100.0\% final-answer pass rate (95\% CI: 96.3--100.0),
30.0\% trace-aware pass rate (95\% CI: 21.9--39.6), and 70 answer-only
missed failures.

\section{Limitations}

This work is an early systems artifact. The benchmark is small and uses
deterministic demo agents for offline reproducibility. The final-answer baseline
is deliberately weak. Offline semantic grading is heuristic. The repair,
learning, fuzzing, and MCP-audit helpers are scaffolding tools rather than
proof-generating systems. The trace schema is versioned, but broader
cross-framework trace compatibility remains future work. These limitations are
documented in the artifact and should be considered when interpreting the
benchmark result.

\section{Related Work}

ReAct studies interleaved reasoning and acting in language-model agents
\cite{yao2023react}. Toolformer studies how language models can learn when to
call tools and how to incorporate tool outputs \cite{schick2023toolformer}.
NeMo Guardrails presents programmable runtime rails for controllable LLM
applications \cite{rebedea2023nemo}. OpenAI's agent-evaluation and tracing
documentation emphasizes trace grading for workflow-level agent issues and
tracing of model calls, tool calls, guardrails, handoffs, and custom events
\cite{openaiAgentEvals,openaiAgentsTracing}. Agent Anvil complements these
directions by packaging local trace assertions, semantic grading, failure
clustering, and CI-compatible artifacts into a small developer workflow.

\section{Conclusion}

Agent Anvil argues for evaluating tool-using agents by their execution traces,
not only by their final answers. The current system is intentionally small:
scenario suites, trace capture, deterministic assertions, optional OpenAI
semantic grading, failure clustering, and CI artifacts. The benchmark result
shows that even a compact trace-aware harness catches workflow failures that a
final-answer baseline misses. Future work should expand the benchmark domains,
add richer assertion semantics, calibrate semantic graders, and provide
adapters for common agent frameworks.

\bibliographystyle{plain}
\bibliography{references}

\end{document}
